\documentclass[12pt]{article}
\usepackage[margin=1in]{geometry}
\usepackage{enumitem}
\usepackage{xcolor}
\usepackage{hyperref}

\pagestyle{empty}

\begin{document}

\begin{center}
{\LARGE\textbf{The Feynman One-Pager}}\\[4pt]
{\large Data Quality Dominates Algorithmic Complexity\\in Neoantigen Immunogenicity Prediction}\\[8pt]
{\small Rao \& Claude Opus 4.6 $\cdot$ March 2026 $\cdot$ \url{github.com/keshav55/cure}}
\end{center}

\vspace{12pt}

\subsection*{The Problem (10 seconds)}

Cancer cells have mutated proteins. Pieces of those proteins get displayed on the cell surface. A cancer vaccine teaches your immune system which pieces to attack. The computational problem: given thousands of pieces, which ones will the immune system actually respond to?

\subsection*{What We Found (30 seconds)}

We tested 125,784 peptide pieces from 131 real cancer patients. Three findings:

\begin{enumerate}[leftmargin=*,itemsep=6pt]

\item \textbf{The \$50 test matters more than any algorithm.}

Every person has 6 HLA alleles---molecular ``display cases'' on their cells. A peptide must fit the patient's specific display case. When we guessed which display cases (3 common ones), our prediction was \textit{worse than random} (AUC = 0.156). When we used the patient's actual display cases, AUC jumped to 0.958.

Same algorithm. Different input. \textbf{The input was the entire bottleneck.}

HLA typing costs \$50--200. It is more valuable than any neural network.

\item \textbf{``How foreign is the mutation'' is the wrong question.}

The most intuitive feature---how chemically different the mutation is from the original amino acid---is \textit{anti-correlated} with success. Immunogenic peptides are \textit{less} foreign on average (AUC = 0.352). We confirmed this across two independent datasets.

Why? Your immune system deletes T-cells that react to highly foreign patterns during development (central tolerance). The most foreign-looking mutations have the fewest remaining T-cells to respond to them.

\item \textbf{Three features and arithmetic beat complex models.}

\vspace{4pt}
\begin{tabular}{ll}
Binding rank (does it fit the display case?) & AUC = 0.966 \\
Stability (how long does it stay displayed?) & AUC = 0.903 \\
Expression (how many copies are made?) & AUC = 0.791 \\
\textbf{All three combined (weighted average)} & \textbf{AUC = 0.958} \\
\end{tabular}
\vspace{4pt}

Published state of the art with neural networks: 0.58--0.81. We beat it with 8 lines of math, because we had the right data.

\end{enumerate}

\subsection*{What We Built}

An open-source pipeline that takes tumor mutations in and produces ranked vaccine candidates with mRNA sequences out. A self-improving daemon (Claude Opus 4.6 + GPT-5.4 Pro) ran 1,048 automated experiments to optimize the scorer. It proposed immunology ideas---conservation-weighted foreignness, TCR repertoire diversity---without being trained in biology.

\subsection*{What It Means}

The bottleneck for personalized cancer vaccines is not algorithms. It is data: which HLA alleles does this patient have, how much is this gene expressed in the tumor, and how stable is the peptide-MHC complex. These are measurable quantities. The first one costs \$50. The pipeline is free.

\vspace{8pt}
\begin{center}
\small\textit{Validated on NeoRanking (M\"{u}ller et al.\ 2023 Immunity), 125,784 peptides, 131 patients.}
\end{center}

\end{document}
